% =============================================================================
% Chapter 3 — Algorithms Used
% =============================================================================
\chapter{Algorithms and Technical Justifications}

\section{Authentication Pipeline}

The authentication flow follows a sequential pipeline in three steps:

\begin{enumerate}
    \item \textbf{Capture} — The webcam captures a 640$\times$480 BGR frame via OpenCV (\texttt{CameraService}).
    \item \textbf{Detection} — The Haar Cascade classifier locates faces in the frame.
    \item \textbf{Recognition} — The ArcFace embedding of the facial ROI is compared to stored templates.
\end{enumerate}

\screenshotplaceholder{Complete facial authentication pipeline}{pipeline-auth}

% ─────────────────────────────────────────────────────────────────────────────
\section{Face Detection — Haar Cascade}

\subsection{Principle}

Detection uses the Viola-Jones cascade classifier (\textbf{Haar Cascade}) provided by OpenCV. This detector scans the image at multiple scales applying Haar features to identify facial characteristics (eyes, nose, mouth).

\subsection{Chosen Parameters}

\begin{table}[H]
    \centering
    \begin{tabular}{lcp{7cm}}
        \toprule
        \textbf{Parameter} & \textbf{Value} & \textbf{Justification} \\
        \midrule
        \texttt{scaleFactor}  & 1.2      & Trade-off between speed and sensitivity. A smaller value (1.05) would be more accurate but slower. \\
        \texttt{minNeighbors} & 6        & Filters false positives: a face must be detected by at least 6 neighboring windows. \\
        \texttt{minSize}      & (80, 80) & Ignores faces that are too small (distant or artifacts). \\
        \bottomrule
    \end{tabular}
    \caption{Haar Cascade detector configuration}
\end{table}

\subsection{Preprocessing}

Before detection, two operations are applied:
\begin{enumerate}
    \item \textbf{Grayscale Conversion} — \texttt{cv2.cvtColor(frame, COLOR\_BGR2GRAY)}
    \item \textbf{Histogram Equalization} — \texttt{cv2.equalizeHist(gray)} to normalize contrast and improve robustness to lighting variations.
\end{enumerate}

\subsection{Choice Justification}

\begin{successbox}[Why Haar Cascade?]
\begin{itemize}[nosep]
    \item \textbf{Speed} — Real-time execution on CPU, suitable for 30 FPS refresh rate.
    \item \textbf{Reliability} — Proven model since 2001, robust under normal lighting conditions.
    \item \textbf{Lightweight} — No GPU dependency, runs on any machine.
    \item \textbf{Complementarity} — Fine-grained recognition is delegated to ArcFace in the second step.
\end{itemize}
\end{successbox}

% ─────────────────────────────────────────────────────────────────────────────
\section{Facial Recognition — ArcFace (InsightFace)}

\subsection{Principle}

ArcFace is a facial recognition model based on a deep convolutional neural network (\textbf{ResNet-50}). It transforms a 112$\times$112 face image into a \textbf{512-dimensional embedding vector}, normalized in L2 norm.

The uniqueness of ArcFace lies in its loss function (\textit{Additive Angular Margin Loss}) which maximizes the angular margin between classes, producing highly discriminative embeddings.

\subsection{Model Used}

\begin{table}[H]
    \centering
    \begin{tabular}{ll}
        \toprule
        \textbf{Property} & \textbf{Value} \\
        \midrule
        Model           & \texttt{w600k\_r50.onnx} (\texttt{buffalo\_l} pack) \\
        Architecture    & ResNet-50 \\
        Input Size      & 112 $\times$ 112 pixels (BGR) \\
        Embedding Size  & 512 dimensions (float32) \\
        Normalization   & L2 (unit vector) \\
        Runtime         & ONNX Runtime (CPU, \texttt{ctx\_id=-1}) \\
        Model Size      & $\sim$280 MB (downloaded once) \\
        \bottomrule
    \end{tabular}
    \caption{Characteristics of the ArcFace model}
\end{table}

\subsection{Embedding Generation Process}

\begin{lstlisting}[language=Python, caption={Building an ArcFace embedding}]
def build_embedding(self, face_roi: np.ndarray) -> np.ndarray:
    # 1. Resize to 112x112
    resized = cv2.resize(face_roi, (112, 112)).astype(np.uint8)
    # 2. Extract features via ArcFace model
    raw_feat = model.get_feat(resized)
    # 3. Flatten and convert to float32
    vector = np.array(raw_feat).flatten().astype(np.float32)
    # 4. Normalize to L2
    norm = np.linalg.norm(vector)
    return vector / norm if norm != 0 else vector
\end{lstlisting}

\subsection{Matching — Euclidean Distance}

The comparison between the captured embedding and the stored templates uses the \textbf{Euclidean distance}:

\begin{equation}
    d(\mathbf{e}_1, \mathbf{e}_2) = \|\mathbf{e}_1 - \mathbf{e}_2\|_2 = \sqrt{\sum_{i=1}^{512} (e_{1,i} - e_{2,i})^2}
\end{equation}

For L2-normalized vectors, the Euclidean distance is related to the cosine similarity:
\begin{equation}
    d^2 = 2(1 - \cos\theta)
\end{equation}

\subsection{Decision Thresholds}

The system uses a \textbf{dual-threshold strategy} to minimize false positives:

\begin{table}[H]
    \centering
    \begin{tabular}{lcl}
        \toprule
        \textbf{Threshold} & \textbf{Value} & \textbf{Decision} \\
        \midrule
        Strict Threshold (\textit{strict\_threshold}) & $\leq 0.75$ & \textcolor{accentgreen}{\textbf{Accepted}} — High confidence \\
        Review Zone (\textit{review\_threshold}) & $0.75 < d \leq 0.90$ & \textcolor{orange}{\textbf{Rejected}} — Insufficient confidence \\
        Beyond & $> 0.90$ & \textcolor{accentred}{\textbf{Unknown}} — No match \\
        \bottomrule
    \end{tabular}
    \caption{Dual-threshold decision strategy}
\end{table}

The distance $\rightarrow$ confidence conversion is: $\text{confidence} = \max(0, \min(1, 1 - d/2))$

% ─────────────────────────────────────────────────────────────────────────────
\section{Digital Watermarking — HMAC-SHA256}

\subsection{Principle}

Each authentication log is protected by a dual integrity mechanism:

\begin{enumerate}
    \item \textbf{SHA-256 Hash} — Payload fingerprint to detect any modification.
    \item \textbf{HMAC-SHA256 Signature} — Cryptographic signature with a secret key to prove authenticity.
\end{enumerate}

\subsection{Payload Construction}

\begin{lstlisting}[language=Python, caption={Watermarking an authentication log}]
payload = "|".join([
    user_name,        # "Alice Martin"
    user_role,        # "Admin"
    status,           # "Authorized"
    f"{confidence:.4f}",  # "0.8523"
    details,          # Matching message
    timestamp.isoformat()  # "2026-05-02T21:30:00"
])
payload_hash = SHA256(payload)           # Fingerprint
signature    = HMAC_SHA256(key, payload) # Watermark
\end{lstlisting}

\subsection{Integrity Verification}

The \textit{Integrity Verification} panel in the admin dashboard recalculates the hashes and signatures of all logs to detect any alteration:

\begin{lstlisting}[language=Python, caption={Log integrity verification}]
computed_hash = SHA256(reconstructed_payload)
hash_ok = (computed_hash == stored_hash)
sign_ok = HMAC_verify(payload, stored_signature)
integrity = hash_ok AND sign_ok
\end{lstlisting}

\subsection{Choice Justification}

\begin{successbox}[Why HMAC-SHA256?]
\begin{itemize}[nosep]
    \item \textbf{Industry Standard} — Used in TLS, JWT, OAuth 2.0.
    \item \textbf{Collision Resistance} — SHA-256 offers $2^{128}$ resistance.
    \item \textbf{Authentication} — The secret key prevents forgery.
    \item \textbf{Performance} — Instant calculation, no impact on UX.
    \item \textbf{Non-alteration} — Any single-bit change invalidates the signature.
\end{itemize}
\end{successbox}

\screenshotplaceholder{Integrity verification panel in the admin dashboard}{panneau-integrite}
